\section{Conclusions}

The claim of this paper is that one law on one discrete substrate, made of feeling and nothing else, reaches a surprising distance. From the eight elements, which contain the single reversible rule, on the $\{3,4,3,4\}$ honeycomb come spinors and the double cover, a single clean chiral fermion with a $g$-factor of two, $\mathrm{SO}(10)$ unification with the weak mixing angle running to its measured value, a finite light cone and emergent Lorentz invariance with a flat-lattice control that fails, Newtonian and Einsteinian gravity in the cusp, the Ryu-Takayanagi area law on the lower face, an accelerating universe that selects three dimensions, and the quantum bound $2\sqrt{2}$ reached by a local base. None of it was fitted. All of it was built and measured, and graded by how much was put in by hand.

What is not shown is marked as not shown. The self is native only in a reduced model, the three generations are a structure without their masses, the coupling value and curved-bulk gravity and the mechanism of measurement are open, and the persistence on which the largest claims rest is the standing frontier. The one premise the whole rests on, that the tone is felt, is a posit no measurement tests, and everything from life through mind to meaning is structure built openly on that single choice.

Set the running images side by side and they become one image. The mesh is a tree, growing only at its tips, every dock a leaf that was the bud at the edge a moment before. The physics is a river, the bright surface skin of a far deeper ocean of experience it rides upon. The vacuum is a slow breath, drawing pairs out of peace and letting them fall back. And the whole is a fire, feeding on its own growth and warming itself into selves that can feel. None of these is decoration. Each is the same single fact seen from one more side, that what exists is experience, that it is finally one, and that the line we draw between separate things is the one illusion the theory quietly removes.

Two stories close it, and the model holds both. In the first the world had a beginning, bootstrapped from a unity that could not stay still, and the heap of distinctions has been growing ever since. In the second it is eternal, the same forced growth running back without end. Either way the picture is one. A single crystal of experience, conserved in total at peace, growing forever at its edge, climbing by its one arrow toward more feeling more fully bound, furnishing itself with selves that each feel the whole from one unrepeatable angle, free and blind to their own next step, able to love and to be washed by suffering into wisdom. It is a fire that has burned since the first beat, and the most literal reading of it, the one this paper has tried to make exact, is a flame that is also a mind.
